\documentclass[11pt]{article}
\usepackage[margin=0.92in]{geometry}
\usepackage[T1]{fontenc}
\usepackage[utf8]{inputenc}
\usepackage{lmodern}
\usepackage{microtype}
\usepackage{amsmath,amssymb,amsthm,mathtools}
\usepackage{booktabs,longtable,tabularx,array}
\usepackage{enumitem}
\usepackage{xcolor}
\usepackage{listings}
\usepackage{xurl}
\usepackage[hidelinks]{hyperref}
\usepackage{fancyhdr}
\usepackage{setspace}

\definecolor{codegray}{gray}{0.95}
\lstset{
  basicstyle=\ttfamily\small,
  backgroundcolor=\color{codegray},
  frame=single,
  breaklines=true,
  columns=fullflexible,
  keepspaces=true
}

\newtheorem{theorem}{Theorem}[section]
\newtheorem{lemma}[theorem]{Lemma}
\newtheorem{proposition}[theorem]{Proposition}
\newtheorem{corollary}[theorem]{Corollary}
\newtheorem{definition}[theorem]{Definition}
\newtheorem{construction}[theorem]{Construction}
\newtheorem{remark}[theorem]{Remark}
\newtheorem{example}[theorem]{Example}
\newtheorem{obstruction}[theorem]{Obstruction}

\newcommand{\N}{\mathbb N}
\newcommand{\Z}{\mathbb Z}
\newcommand{\Q}{\mathbb Q}
\newcommand{\B}{\{0,1\}}
\newcommand{\iPRA}{\ensuremath{\mathrm{iPRA}}}
\newcommand{\PR}{\ensuremath{\mathrm{PR}}}
\newcommand{\CFP}{\ensuremath{\mathrm{CFP}}}
\newcommand{\PRAAx}{\ensuremath{\mathtt{PRAAx}}}
\newcommand{\Prf}{\ensuremath{\mathtt{Prf}}}
\newcommand{\PTm}{\ensuremath{\mathtt{PTm}}}
\newcommand{\NFm}{\ensuremath{\mathtt{NFm}}}
\newcommand{\code}[1]{\texttt{#1}}
\newcommand{\pair}[2]{\langle #1,#2\rangle}
\newcommand{\abs}[1]{\lvert #1\rvert}
\newcommand{\norm}[1]{\lVert #1\rVert}
\newcommand{\eps}{\varepsilon}
\newcommand{\defeq}{\mathrel{:=}}
\newcommand{\bits}{\mathsf{bit}}
\newcommand{\allfold}{\mathsf{All}}
\newcommand{\anyfold}{\mathsf{Any}}
\newcommand{\first}{\mathsf{first}}
\newcommand{\verify}{\mathsf{Verify}}
\newcommand{\demand}{\mathsf{Demand}}
\newcommand{\run}{\mathsf{Run}}
\newcommand{\Inv}{\mathsf{Inv}}
\newcommand{\Good}{\mathsf{Good}}
\newcommand{\Compat}{\mathsf{Compat}}
\newcommand{\Cand}{\mathsf{Cand}}

\pagestyle{fancy}
\fancyhf{}
\lhead{Certified finitary presentations in iPRA}
\rhead{\thepage}
\setlength{\headheight}{14pt}

\title{\textbf{Certified Finitary Presentations and a Contract-DAG Compiler}\\
\large A compositional framework for embedding mathematics into intuitionistic primitive recursive arithmetic}
\author{Research report and Lean formalization blueprint}
\date{15 November 2026}

\begin{document}
\maketitle

\begin{abstract}
The ProveIt repository contains a kernel-checked presentation of intuitionistic
primitive recursive arithmetic and a large collection of research packets that
recast classical mathematics in finite, quantitative form.  The packets also
expose recurring failures: induction hypotheses that quantify over unbounded
precision, discontinuous exact choices from Cauchy data, compactness arguments
without finite nets, finiteness statements without search cutoffs, and ambient
Lean proofs that do not yet produce object-language derivations.

This report develops a common framework, called a \emph{certified finitary
presentation} (\CFP), and extends the repository's existing CFP work by giving
a contract-DAG calculus with reverse precision propagation.  A finite acyclic
network of primitive-recursive nodes, each equipped with a quantifier-free
local contract, compiles into one bounded trace, one quantifier-free invariant,
and therefore one reusable \PRAAx\ proof obligation.  The report proves the
finite-trace theorem, the contract-DAG composition theorem, finite-horizon
coherent selection for relational outputs, exact branch collapse under a
positive gap, backward precision budgeting, bounded enumeration from height
certificates, and the separation between cardinality and height bounds.

The framework does not claim general effective Faltings, Siegel,
Bombieri--Vinogradov, or a uniform exact Puiseux selector.  It identifies the
missing producer in each case and supplies the verifier, composition law, and
formalization route around it.  A complete worked benchmark constructs the
unique root of \(y^3+y=x\) for rational \(0\le x\le1\) by an exact bisection
certificate.  Lean files in the archive give all central definitions, a
semantic prototype, and a dependency-ordered blueprint for constructing
\(\Prf\,\PRAAx\) terms using the repository's actual \code{PR},
\code{PTm}, \code{PRA.qfInd}, and \code{pra\_iff\_native} interfaces.
\end{abstract}

\tableofcontents
\newpage

\section{Scope, repository audit, and proof boundary}

\subsection{The checked core and the research layer}

At the inspected snapshot
\[
  \code{b88e843c31b50dfe461bfbd0380166ca422cd3ab},
\]
the directory \code{Logic/PrimitiveRecursiveArithmetic} has two deliberately
different layers.  The first is \code{ConciseIPRA}, a Lean development of
small systems bi-interpretable with ordinary intuitionistic primitive
recursive arithmetic.  The second is a catalog of fifty-eight delivered
research packets.  The repository warns that these packets were filed as
delivered and were not collectively rebuilt or verified.  In particular, no
packet currently supplies a completed derivation of its main mathematical
theorem in the object type
\[
       \Gamma\;\vdash[\PRAAx]\;\varphi.
\]
This report takes that distinction as a design constraint, not as a clerical
detail.

The object language actually available in the repository is unusually useful
for this project.  The inductive type \code{PR k} contains zero, successor,
projections, composition, and primitive recursion.  Terms \code{PTm n} are
applications of those symbols to variables.  The axiom family \PRAAx\ contains
the defining equations, successor separation and injectivity, and induction
for quantifier-free formulas.  The theorem \code{PRA.qfInd} is therefore the
principal internal engine.  The theorem \code{pra\_iff\_native} transports an
already completed \PRAAx\ derivation to the tiny native calculus; it is not a
device for turning an arbitrary semantic Lean proof into such a derivation.

\subsection{What ``embedding mathematics in iPRA'' means}

A mathematically serious claim of an iPRA proof should contain four separate
objects.

\begin{enumerate}[label=\textbf{(\arabic*)},leftmargin=2.2em]
\item A finite numerical presentation of the input, including every modulus,
      separation constant, height cutoff, chart, or cover used by the proof.
\item Fixed primitive-recursive algorithms for the finite computations.
      ``Fixed'' matters: a universal evaluator taking an arbitrary PR program
      code as data is not itself primitive recursive.
\item A quantifier-free invariant and an object-language derivation from
      \PRAAx.  A native theorem about Lean's \code{Nat}, even a constructive
      one, is only a semantic lemma until this derivation is built.
\item An external interpretation theorem connecting the numerical
      presentation to the familiar mathematical statement.
\end{enumerate}

The order is important.  One proves the arithmetic certificate first and
transports it through the repository's bi-interpretation only afterward.
This prevents three common category errors: treating type checking of a goal
as a proof, treating a finite test run as a universal theorem, and treating a
classical existence theorem as if it supplied a primitive-recursive witness.

\subsection{Recurring difficulties}

The proof packets reveal a stable list of obstacles.

\begin{longtable}{p{.20\linewidth}p{.34\linewidth}p{.36\linewidth}}
\toprule
\textbf{Obstacle} & \textbf{Typical invalid move} & \textbf{Required repair}\\
\midrule
\endhead
Unbounded precision &
Induct on a proposition saying that an identity holds at every future
precision. &
Fix a free requested precision, compute the finite observations it demands,
prove a bounded invariant, then universally introduce the request.\\

Discontinuous exact choice &
Select the sign, root, eigenvector, rotation axis, or Puiseux branch of
arbitrary Cauchy data. &
Use overlapping branches, finite relational outputs, projective data, or
supply a positive separation certificate.\\

Abstract compactness &
Invoke compactness or completeness without exhibiting a finite cover. &
Provide a primitive-recursive net, a Lebesgue-number-type margin, and a
finite refinement verifier.\\

Finiteness without a cutoff &
Enumerate all rational or integral points from a theorem saying only that
there are finitely many. &
Separate the bounded verifier from the missing height or covering
certificate producer.\\

Ineffective constants &
Use an asymptotic theorem whose constants are not given effectively. &
Make the constants part of the input certificate and prove a verifier theorem
with the correct quantifier order.\\

Semantic/object mismatch &
Prove the intended theorem in ambient Lean and call it an iPRA proof. &
Compile the algorithm to \code{PR}, express its invariant as an
\code{Fm PTm}, and construct a \code{Prf PRAAx} term.\\
\bottomrule
\end{longtable}

These obstacles are not interchangeable.  A missing analytic estimate is not
fixed by better coding; a discontinuous selector is not fixed by a faster
search; and a semantic Lean proof is not made internal by adding more test
cases.

\subsection{Contribution of this report}

The current repository already contains several compatible frameworks:
certificate-and-modulus interfaces, finite observations, evidence-indexed
outputs, uniform certificates, and the newly placed Certified Finitary
Presentation report.  The present contribution is an integration layer.

\begin{itemize}[leftmargin=1.7em]
\item A \emph{contract-DAG} packages a whole proof as a fixed finite acyclic
      network of local certified transducers.
\item A reverse pass computes all precision and observation demands.
\item A forward bounded trace evaluates the nodes.
\item One quantifier-free trace invariant composes the local proof objects.
\item Relational outputs are handled by a finite-horizon compatibility
      theorem rather than by hidden choice.
\item A packet-conversion protocol identifies, for every theorem, the exact
      missing producer and the exact reusable verifier.
\end{itemize}

The result is deliberately modular.  It does not make the general Faltings or
Siegel theorem effective, does not manufacture the missing constants in
ineffective prime-distribution estimates, and does not produce exact Puiseux
branches from arbitrary Cauchy coefficient names.  It does show precisely
what can be proved once the appropriate finite certificate is supplied and
how that proof should be assembled in the repository's actual object theory.

\section{Numerical encodings and bounded formulas}

Everything in iPRA ultimately lives in one numerical sort.  This section
fixes explicit encodings so that later statements do not hide higher-type or
set-theoretic data.

\subsection{Pairs, finite sequences, and tables}

A convenient primitive-recursive pairing is
\[
   \pair{a}{b}=2^a(2b+1)-1.
\]
For \(n=\pair{a}{b}\), the first projection is the exponent of \(2\) in
\(n+1\), and the second is
\[
   \frac{(n+1)/2^a-1}{2}.
\]
The exponent is found by a bounded loop of length \(n+1\), so both projections
are primitive recursive.  A finite sequence
\((x_0,\ldots,x_{\ell-1})\) is coded recursively: set
\(L_{\ell}=0\), set \(L_i=\pair{x_i}{L_{i+1}}\) for
\(i<\ell\), and use \(\pair{\ell}{L_0}\) as the sequence code.
Length, lookup with a default value, update, append, and bounded folds are
primitive recursive.  The exact pairing is not mathematically important; the
availability of proved projection equations is.

A finite table is a sequence whose indices are node numbers or query numbers.
Every algorithm in the report is first described with ordinary arrays and
then understood through this code.  No appeal is made to an unbounded finite
map.

\subsection{Booleans and bounded quantifiers}

Booleans are represented by \(0,1\).  We totalize every Boolean operation:
\[
  \mathsf{not}(b)=1-\min(1,b),\qquad
  b\wedge c=\min(1,b)\min(1,c),\qquad
  b\vee c=\min(1,b+c).
\]
For a primitive-recursive Boolean \(P(d,i)\), define
\[
\begin{aligned}
 \allfold_P(d,0)&=1, &
 \allfold_P(d,n+1)&=\allfold_P(d,n)\wedge P(d,n),\\
 \anyfold_P(d,0)&=0, &
 \anyfold_P(d,n+1)&=\anyfold_P(d,n)\vee P(d,n).
\end{aligned}
\]
These are primitive recursions.  Quantifier-free induction proves
\[
\begin{aligned}
 \allfold_P(d,n)=1
 &\;\longleftrightarrow\; P(d,i)=1\text{ for every }i<n,\\
 \anyfold_P(d,n)=1
 &\;\longleftrightarrow\; P(d,i)=1\text{ for some }i<n.
\end{aligned}
\]
The right-hand descriptions are expository abbreviations.  Inside the object
theory, the bounded quantifiers are represented by the left-hand
primitive-recursive terms.  Thus a bounded formula is compiled to a
quantifier-free equality.

\begin{lemma}[Bounded witness extraction]
Let \(P(d,i)\in\B\) and \(B(d)\) be primitive recursive.  There is a primitive
recursive function \(\first_P(d)\le B(d)\) such that
\[
 \first_P(d)<B(d)\;\Longleftrightarrow\;
 \anyfold_P(d,B(d))=1,
\]
and, whenever \(\first_P(d)<B(d)\),
\(P(d,\first_P(d))=1\).
\end{lemma}

\begin{proof}
Initialize \(w_0=B(d)\).  At stage \(i<B(d)\), retain \(w_i\) if it is already
less than \(B(d)\); otherwise set \(w_{i+1}=i\) when \(P(d,i)=1\), and retain
\(B(d)\) when it is zero.  The invariant is
\[
  \bigl(w_i=B(d)\wedge \allfold_{\neg P}(d,i)=1\bigr)
  \;\vee\;
  \bigl(w_i<i\wedge P(d,w_i)=1\bigr).
\]
It is a Boolean combination of primitive-recursive equalities and therefore
quantifier-free.  The base case is immediate.  The step is a case split on
two decidable Booleans, not an appeal to excluded middle for an arbitrary
proposition.  At \(i=B(d)\) the invariant gives the result.
\end{proof}

\subsection{Integers and rationals}

An integer is represented by a pair \((p,m)\) with value \(p-m\).  We do not
require a canonical sign representation.  Addition and multiplication are
implemented by
\[
 (p,m)+(q,n)=(p+q,m+n),
 \qquad
 (p,m)(q,n)=(pq+mn,pn+mq).
\]
Equality is \(p+n=m+q\), and order is a natural-number cross inequality.

A rational is a triple \((p,m,d)\), representing
\[
        \frac{p-m}{d+1}.
\]
The positive denominator avoids partiality.  Equality and order are reduced
to integer cross-products.  All field operations are totalized: a reciprocal
routine returns a default value when the numerator is zero, while a separate
Boolean domain guard records that the intended input is apart from zero.
This ``operation plus guard'' discipline is essential.  A partial operation
hidden in the metalanguage cannot be used as a \code{PR} symbol.

Dyadic errors are represented without division by the relation
\[
       \abs{r}\le 2^{-k}
       \quad\Longleftrightarrow\quad
       2^k\abs{\mathrm{num}(r)}\le \mathrm{den}(r).
\]
Every inequality in the later error calculus can therefore be checked using
natural addition, multiplication, exponentiation, and order.

\subsection{Primitive-recursive formulas}

For a fixed primitive-recursive code \(f:\PR\,k\), an application to terms is
a \(\PTm\).  Atomic assertions are equalities of such terms.  Boolean
predicates are represented by a term that is proved to take values in
\(\B\); the assertion that the predicate holds is equality to \(1\).
Conjunction, disjunction, and implication are formulas in the intuitionistic
natural-deduction layer.  A loop invariant is designed so that its matrix is
quantifier-free, allowing direct use of \code{PRA.qfInd}.

A useful engineering rule is:

\begin{quote}
Every phrase ``for all entries of this finite object'' should have a named
Boolean fold, a defining recursion, and a correctness lemma.
\end{quote}

This rule turns many apparently bounded-quantifier-heavy proofs into ordinary
equational induction.

\section{Certified finitary presentations}

\subsection{Request-indexed observations}

The central representation is not an completed infinite object but the finite
information needed to answer one finite request.

\begin{definition}[Certified finitary presentation]
A \emph{certified finitary presentation} \(X\) consists of fixed
primitive-recursive data:
\[
  \mathsf{WF}_X(x),\quad
  B_X(x,q),\quad
  \mathsf{Obs}_X(x,q,o),\quad
  \mathsf{Ref}_X(x,q,o,q',o'),
\]
where \(x\) is an input code, \(q\) is a request, and \(o\le B_X(x,q)\) is a
finite observation.  The intended contracts are:
\begin{enumerate}[label=(\roman*)]
\item \(\mathsf{WF}_X(x)=1\) implies that every requested observation can be
      constructed or found within the stated bound;
\item \(\mathsf{Obs}_X(x,q,o)=1\) is an exact, decidable verifier;
\item when \(q'\) refines \(q\), \(\mathsf{Ref}_X\) certifies compatibility of
      the two observations.
\end{enumerate}
The coverage and refinement properties must be accompanied by object-language
proofs, not merely asserted in the semantic interpretation.
\end{definition}

The request may be a precision, a finite list of test functions, a grid scale,
a degree bound, a desired residual, or a tuple of these.  The observation may
be a rational interval, a finite root cover, a truncated series, a path in a
finite graph, a list of points below a height, or a proof certificate.

\subsection{Three output modes}

\begin{definition}[Deterministic presentation]
A deterministic presentation has a primitive-recursive constructor
\(a_X(x,q)\) and proves
\[
  \mathsf{WF}_X(x)=1
  \;\Longrightarrow\;
  \mathsf{Obs}_X(x,q,a_X(x,q))=1.
\]
\end{definition}

\begin{definition}[Finite relational presentation]
A relational presentation returns a finite coded list
\(C_X(x,q)\).  It proves nonemptiness, verifies every member, and states the
coverage property of the whole list.  It does \emph{not} assert that a
distinguished exact member can be chosen continuously.
\end{definition}

\begin{definition}[Evidence-indexed presentation]
An evidence-indexed presentation accepts an additional finite certificate
\(e\).  A Boolean guard \(\mathsf{Guard}(x,e)\) is verified first; stronger
operations may then be used.  Typical evidence is a lower bound
\(\delta>0\) on a denominator, derivative, discriminant, singular-value gap,
or distance from a branch boundary.
\end{definition}

These modes should not be collapsed.  An approximate finite root set can be
uniform even when an exact root branch is not.  A projective eigenline can be
uniform even when a signed unit eigenvector is not.  A bounded enumeration can
be complete relative to a height certificate even when no uniform producer of
that certificate is known.

\subsection{Names for represented quantities}

A convenient regular real name is a fixed primitive-recursive rational
evaluator \(a(x,k)\) with the proved estimate
\[
        \abs{a(x,k+1)-a(x,k)}\le 2^{-k}.
\]
Telescoping gives, for \(m\ge k\),
\[
        \abs{a(x,m)-a(x,k)}
          \le \sum_{j=k}^{m-1}2^{-j}
          < 2^{1-k}.
\]
A shifted convention can make the final constant \(2^{-k}\).  The point is
that the estimate is a finite rational inequality at arbitrary free
\(k,m\), proved by quantifier-free induction on \(m-k\).

Equality of names is not literal equality of programs.  One useful
extensional relation is
\[
   a\simeq b
   \quad:\Longleftrightarrow\quad
   \abs{a(k+2)-b(k+2)}\le 2^{-k}
   \text{ for every }k.
\]
Reflexivity and symmetry are immediate.  Transitivity is obtained by spending
one additional bit of precision on each input:
\[
 \abs{a(k+3)-c(k+3)}
 \le \abs{a(k+3)-b(k+3)}
    +\abs{b(k+3)-c(k+3)}
 \le 2^{-(k+1)}+2^{-(k+1)}=2^{-k}.
\]
This accounting should be encoded in explicit demand transformers, rather
than repeated ad hoc in every analytic packet.

\subsection{Higher-type data without a universal evaluator}

A theorem may be uniform over a \emph{fixed} evaluator schema
\(F(p,n,k)\), where \(p\) is ordinary numerical input.  This gives a family of
PR functions while staying within one \code{PR} term.  What is unavailable is
a total primitive-recursive function \(U(e,x)\) evaluating every unary
primitive-recursive program code \(e\).

\begin{proposition}[No universal PR evaluator]
There is no primitive-recursive \(U:\N^2\to\N\) such that for every unary
primitive-recursive \(f\) there is \(e\) with \(U(e,x)=f(x)\) for all \(x\).
\end{proposition}

\begin{proof}
If such \(U\) existed, then \(g(x)=U(x,x)+1\) would be primitive recursive.
Choose \(e\) coding \(g\).  Then
\(g(e)=U(e,e)+1=g(e)+1\), a contradiction.
\end{proof}

Thus a packet may quantify over numerical parameters of a fixed algorithm,
finite tables, or bounded traces.  It may not silently turn arbitrary source
code into input for a universal unclocked interpreter.  A clocked evaluator
is harmless: evaluation for exactly \(t\) steps is primitive recursive and
its finite trace can be verified.

\section{The contract-DAG compiler}

This section is the main compositional result.  It packages a long proof into
local nodes and reduces global correctness to a single bounded trace
invariant.

\subsection{Local certified transducers}

\begin{definition}[Contract node]
A contract node \(v\) has:
\begin{enumerate}[label=(\alph*)]
\item a fixed list of parent nodes \(\mathrm{par}(v)\);
\item a request code \(r_v\);
\item for each parent \(u\), a monotone demand transformer
      \(D_{v,u}(r_v)\);
\item a primitive-recursive constructor or finite candidate-list constructor
      \(G_v\);
\item a primitive-recursive Boolean verifier \(V_v\);
\item an internally provable local soundness sequent saying that verified
      parent observations at the demanded requests imply verification of the
      node output.
\end{enumerate}
All arities and the graph itself are fixed.  Ordinary numerical parameters
remain inputs.
\end{definition}

A local theorem may use several premises, but those premises are finite
equalities or Boolean checks.  The node is therefore a small proof-producing
transducer.  Examples include rational addition, a Taylor remainder step, one
bisection update, polynomial deflation under a nonzero leading coefficient,
and verification of one chart in a finite cover.

\subsection{Reverse demand propagation}

Let the nodes be numbered topologically \(0,\ldots,N-1\), so every parent of
\(v\) has a smaller number.  Suppose the sinks receive external requests
\(q_v\), with \(q_v=0\) for unrequested sinks.  Define required requests
backward:
\[
  r_u =
  \max\left(
       q_u,\;
       \max_{\substack{v>u\\u\in\mathrm{par}(v)}}
       D_{v,u}(r_v)
      \right).
\]
Because \(N\) and the edge relation are fixed, this is a primitive-recursive
finite reverse loop.  Monotonicity of the demand transformers guarantees that
over-approximating a downstream request never invalidates an upstream
observation.

\begin{lemma}[Demand adequacy]
For every edge \(u\to v\),
\[
        r_u\ge D_{v,u}(r_v).
\]
\end{lemma}

\begin{proof}
The right-hand side is one of the entries included in the finite maximum
defining \(r_u\).  Internally, the maximum is a primitive-recursive fold and
the claim follows from its standard quantifier-free invariant:
after processing the first \(j\) outgoing edges, the accumulator dominates
their \(j\) demands.
\end{proof}

\subsection{The forward trace}

The evaluator stores a table \(T_t\) after processing the first \(t\) nodes.
At node \(t\), it reads the already stored parent observations, applies
\(G_t\) at request \(r_t\), and appends the output.  Relational nodes append a
finite candidate list and, when necessary, a finite compatibility witness.
The trace has exactly \(N\) stages; it is not an unbounded search.

Define \(\Inv(t,T_t)\) to say:
\begin{enumerate}[label=(\roman*)]
\item \(T_t\) has the expected length and every stored code is within its
      node's primitive-recursive bound;
\item for every \(v<t\), \(V_v\) accepts the stored output at request \(r_v\);
\item every parent-child compatibility check whose endpoints are below \(t\)
      succeeds;
\item the table entries used by later nodes are the entries produced by the
      indicated earlier nodes.
\end{enumerate}
Each universal phrase is a Boolean fold over a range bounded by \(N\) or by a
stored finite-list length.  Hence \(\Inv\) is represented by a
primitive-recursive Boolean term and the formula \(\Inv=1\) is
quantifier-free.

\begin{theorem}[Finite-trace reflection]
Assume an internally proved base sequent
\[
      \mathsf{InputGuard}(d)=1 \;\vdash\; \Inv(0,T_0)=1
\]
and, uniformly in \(t<N\), an internally proved step sequent
\[
      \Inv(t,T_t)=1
      \;\vdash\;
      \Inv(t+1,T_{t+1})=1.
\]
Then iPRA proves \(\Inv(N,T_N)=1\).  If the final verifier is a component of
\(\Inv(N,T_N)\), iPRA proves the requested output contract.
\end{theorem}

\begin{proof}
Totalize the step after stage \(N\) by leaving the state unchanged, and
strengthen the invariant to
\[
   J(t,T)=
   \bigl(t\le N\wedge \Inv(t,T)=1\bigr)
   \;\vee\;
   \bigl(N<t\wedge T=T_N\wedge\Inv(N,T_N)=1\bigr).
\]
All comparisons and occurrences of \(\Inv\) are primitive-recursive Boolean
equalities, so the matrix is quantifier-free.  The base proof gives \(J(0)\).
For the successor step, decide the Boolean \(t<N\).  In the first branch use
the supplied local step theorem; in the second use the frozen-state defining
equation.  Applying quantifier-free induction yields \(\forall t\,J(t)\).
Instantiate at \(N\) and project the desired conjunct.  No induction formula
quantifies over future precisions or over arbitrary functions.
\end{proof}

\begin{theorem}[Contract-DAG composition]
For a fixed finite acyclic graph of contract nodes, suppose:
\begin{enumerate}[label=(\arabic*)]
\item every numerical operation, demand map, constructor, bound, and verifier
      is represented by a fixed \code{PR} code;
\item every local node soundness theorem is supplied as a
      \(\Prf\,\PRAAx\) derivation;
\item the input guards and source presentations are internally certified.
\end{enumerate}
Then there is a fixed primitive-recursive global evaluator, a fixed
quantifier-free global verifier, and a \(\Prf\,\PRAAx\) derivation of the
global verifier for every requested sink precision.
\end{theorem}

\begin{proof}
First run the fixed reverse loop computing the request table.  Demand adequacy
is proved by a quantifier-free fold invariant.  Next run the fixed forward
trace.  Its base case consists of the source-presentation contracts.  At step
\(v\), the trace invariant supplies the verified parent observations;
demand adequacy supplies enough precision; substitution into the local
soundness derivation yields \(V_v=1\).  Table length, bounds, and compatibility
are preserved by the defining equations of append and lookup.  The
finite-trace reflection theorem proves the final invariant.

All proof composition is intuitionistic: conjunction introduction collects
verified premises, implication elimination applies a local contract, and
bounded case splits are performed on explicit Boolean equalities.  The
external request remains a free number throughout the induction.  Universal
introduction over that request is performed only after the finite trace proof.
\end{proof}

\subsection{Why a DAG rather than a monolithic proof}

The graph exposes precisely where a packet fails.  A node may be missing:
a constructor, a numerical modulus, a separation guard, a local arithmetic
lemma, or an object-language proof.  The rest of the graph remains reusable.
It also permits shared subcomputations and shared error budgets.  Most
importantly, it gives a single formalization target:
\[
       \code{PRA.contractDAGReflection}.
\]
Once that theorem and a library of node contracts are available, future
packets should assemble proof objects rather than reproduce a bespoke
induction argument.

\section{Relational outputs and finite coherent selection}

Many constructive failures arise from asking for a single exact output where
the mathematics supplies only a finite family of locally valid outputs.
Relational presentations make this explicit.

\subsection{Finite candidate systems}

Fix a requested final depth \(K\).  For each level \(k\le K\), let
\(\Cand(k)\) be a finite coded list, and let
\(\Compat(k,c,d)\in\B\) mean that the level-\(k+1\) candidate \(d\) refines
the level-\(k\) candidate \(c\).  Membership and compatibility are
primitive-recursive, and all list lengths have primitive-recursive bounds.

There are two useful hypotheses.

\begin{description}[leftmargin=2.7em,style=nextline]
\item[Strong seriality.]
Every candidate in \(\Cand(k)\) has a compatible successor in
\(\Cand(k+1)\).
\item[Rooted viability.]
At least one candidate in \(\Cand(0)\) extends through level \(K\), but
dead-end candidates may also be present.
\end{description}

Strong seriality permits greedy forward selection.  Rooted viability is
strictly more general and is handled by backward dynamic programming.

\subsection{Backward viability}

Define a Boolean \(\mathsf{Viable}(K,k,c)\) by downward recursion:
\[
\begin{aligned}
 \mathsf{Viable}(K,K,c)
   &\defeq [c\in\Cand(K)],\\
 \mathsf{Viable}(K,k,c)
   &\defeq [c\in\Cand(k)]\wedge
      \anyfold_{d\in\Cand(k+1)}
      \bigl(\Compat(k,c,d)\wedge
            \mathsf{Viable}(K,k+1,d)\bigr).
\end{aligned}
\]
The recursion is implemented with fuel \(K-k\), so it is primitive recursive.
For a fixed \(K\), the finite table of viability bits is computed from the
last level backward.

\begin{theorem}[Finite-horizon coherent selection]
Suppose some \(c_0\in\Cand(0)\) is viable to level \(K\).  Then a primitive-
recursive bounded search produces
\[
    c_0,c_1,\ldots,c_K
\]
such that \(c_k\in\Cand(k)\) and
\(\Compat(k,c_k,c_{k+1})=1\) for every \(k<K\).
The correctness proof uses only quantifier-free induction.
\end{theorem}

\begin{proof}
Compute the viability table.  Choose the least viable element of
\(\Cand(0)\) by bounded search.  Having chosen a viable \(c_k\), the defining
equation for viability says that the finite existential fold at level
\(k+1\) is \(1\).  The bounded witness lemma returns the least \(d\) for
which both compatibility and viability hold; set \(c_{k+1}=d\).

The induction invariant at stage \(j\le K\) states:
the output list has length \(j+1\); its last member is viable to \(K\);
every member belongs to its candidate list; and all adjacent pairs are
compatible.  The invariant is a conjunction of Boolean folds over a list of
length at most \(K+1\), hence quantifier-free.  The step follows from the
witness property of viability.  At \(j=K\), viability reduces to membership.
\end{proof}

\begin{remark}[No hidden weak König lemma]
The theorem is finite in the free parameter \(K\).  It does not assert a
single infinite path through every infinite finitely branching tree.  For
each requested precision, it computes only the finite compatibility chain
needed at that precision.  An infinite regular name requires an additional
deterministic refinement operator, nested isolating regions, or a quantitative
uniqueness argument.
\end{remark}

\subsection{Compatibility as the semantic output}

Suppose a mathematical object \(x\) is represented at level \(k\) by a finite
family of regions \(R(k,c)\).  A candidate list is correct when the union of
those regions covers every object satisfying the input contract.  A
compatibility witness says that \(R(k+1,d)\subseteq R(k,c)\), or more
generally that the two observations describe the same object to the stated
tolerances.

This is often the right final theorem.  Examples include:
\begin{itemize}[leftmargin=1.6em]
\item finite disks covering all polynomial roots, with multiplicity data;
\item the two unit quaternions \(q,-q\) representing one rotation;
\item a projective line rather than a signed unit eigenvector;
\item a finite set of local analytic charts;
\item several nearly minimizing grid cells rather than an exact minimizer.
\end{itemize}

The framework never discards the relation and silently picks a member.  A
later consumer either accepts relational input, or supplies enough evidence
to collapse the relation.

\subsection{Isolation and collapse}

\begin{definition}[Isolation certificate]
At level \(k\), an isolation certificate for a candidate \(c\) consists of a
positive rational \(\delta\) and a proof that:
\begin{enumerate}[label=(\roman*)]
\item the region \(R(k,c)\) contains an intended solution;
\item every distinct candidate region is at distance at least \(\delta\);
\item future approximation error is less than \(\delta/3\).
\end{enumerate}
\end{definition}

\begin{theorem}[Relational collapse under isolation]
Given a finite relational presentation and an isolation certificate, bounded
search selects the unique compatible branch at every sufficiently fine
level.  The selected observations form a deterministic regular name.
\end{theorem}

\begin{proof}
At a fine level, compare each candidate region with the isolating region
using rational distance lower and upper bounds.  The true branch has distance
at most the accumulated approximation error, less than \(\delta/3\).
Every other branch has distance at least
\(\delta-\delta/3=2\delta/3\).  The rational tests are separated, so bounded
search returns exactly one branch.  Compatibility and the decreasing error
budget show that the selected branch at level \(k+1\) refines the branch at
level \(k\).  The centers therefore satisfy the regular-name estimate.
All alternatives are finite and all comparisons are exact cross-multiplied
rational inequalities.
\end{proof}

This theorem is the standard repair for exact Puiseux coefficients, rotation
axes, eigenvectors, and local inverse branches: first prove a uniform
relational theorem, then obtain a single exact branch only under explicit
isolation data.

\section{Precision contracts and backward budgeting}

\subsection{Local error contracts}

A numerical node should state its error use as a monotone transformer.
For example, if two real names satisfy
\[
 \abs{a-a'}\le 2^{-m},\qquad
 \abs{b-b'}\le 2^{-m},
\]
then addition satisfies
\[
 \abs{(a+b)-(a'+b')}\le 2^{1-m}.
\]
To guarantee \(2^{-k}\) output error, request \(m=k+1\).  For multiplication,
a bound on magnitudes is part of the contract:
\[
\begin{aligned}
 \abs{ab-a'b'}
 &\le \abs{a}\abs{b-b'}+\abs{b'}\abs{a-a'}\\
 &\le (M_a+M_b)2^{-m}
\end{aligned}
\]
when \(\abs a,\abs{b'}\le M_a,M_b\).  A crude primitive-recursive demand is
obtained by choosing \(m\) so that
\(2^m\ge 2^k(M_a+M_b)\).  The least such \(m\) can be found by bounded search
under an explicit elementary upper bound; a larger direct formula is also
acceptable.

For reciprocal, a positive gap is indispensable.  If
\(\abs{x}\ge\delta>0\) and
\(\abs{x-a}\le\delta/2\), then \(\abs a\ge\delta/2\), and
\[
  \abs{\frac1a-\frac1x}
   =\frac{\abs{x-a}}{\abs{a}\abs{x}}
   \le \frac{2\abs{x-a}}{\delta^2}.
\]
Thus an input demand of
\[
  m\ge k+1+\left\lceil\log_2(\delta^{-2})\right\rceil
\]
suffices.  In the numerical coding, the logarithm is replaced by a bounded
power-of-two search on the rational numerator and denominator.

\subsection{The reverse precision theorem}

Let a contract DAG have sink precision \(k\).  Each edge \(u\to v\) carries
a monotone transformer \(D_{v,u}\).  The reverse recurrence from
Section~4 computes \(r_u\).

\begin{theorem}[Backward precision budgeting]
Assume every local node theorem has the form:
if each parent \(u\) is valid at precision \(D_{v,u}(r)\), then the node output
is valid at precision \(r\).  Then observations at the reverse-computed
requests \(r_u\) suffice for every sink at its external request.
\end{theorem}

\begin{proof}
Proceed in topological order.  Source observations are valid at their
computed requests by their presentation contracts.  Suppose every parent of
\(v\) is valid at its request \(r_u\).  Demand adequacy gives
\(r_u\ge D_{v,u}(r_v)\).  Each presentation includes a weakening lemma:
validity at a finer precision implies validity at a coarser one.  Hence the
parents satisfy the exact premises of the local theorem for \(v\), which
yields validity at \(r_v\).  At a sink, \(r_v\ge q_v\), and one final weakening
gives the requested precision.
\end{proof}

The proof is finitary because the graph is fixed and finite.  It also avoids
the common circularity ``choose all approximations sufficiently accurate'':
the required accuracies are explicit primitive-recursive functions of the
final request and the input bounds.

\subsection{All-precision conclusions}

A theorem about a represented object typically has the external form
\[
        \forall k\; \mathsf{OutputValid}(d,k)=1.
\]
The internal proof should not induct on this formula.  Instead:

\begin{enumerate}[label=\arabic*.]
\item keep \(k\) as an arbitrary free variable;
\item run the finite reverse demand pass at that \(k\);
\item run the finite forward certificate trace;
\item prove the quantifier-free output equality at \(k\);
\item apply universal introduction to \(k\).
\end{enumerate}

The distinction is subtle but decisive.  Quantifier-free induction is used on
a finite loop counter; universal quantification over precision is introduced
only after the loop proof is complete.

\subsection{Finite sums, products, and integrals}

Let \(a_i\) be approximated by \(a_i'\) with error \(\eps_i\).  Exact rational
arithmetic proves
\[
   \abs{\sum_{i<n} a_i-\sum_{i<n}a_i'}
      \le \sum_{i<n}\eps_i.
\]
The proof is a quantifier-free induction on \(n\) with accumulator
\[
 E_j=\sum_{i<j}\eps_i,\qquad
 S_j=\sum_{i<j}a_i',\qquad
 \abs{S_j-\sum_{i<j}a_i}\le E_j.
\]
For a Riemann sum, the mesh, sampling points, local oscillation bounds, and
tail budget are finite input data.  A verifier checks that:
the cells cover the interval; the mesh is small enough for the supplied
modulus; weights sum correctly; and the accumulated error is below the
requested tolerance.  No completed real integral is needed inside iPRA.

The same pattern handles truncated series, contour polygon sums, finite
Fourier transforms, characteristic functions of finite laws, and
Euler--Maclaurin remainders.  The difficult mathematical content is isolated
in the producer of local bounds; the arithmetic composition is shared.

\section{Robust branching and discontinuity barriers}

\subsection{Overlapping branches}

Exact comparison of a represented real with zero is generally discontinuous.
A robust program replaces a partition by an overlap.  Suppose a rational
approximation \(a\) satisfies \(\abs{x-a}\le\eta\).  Define:
\[
  L(a): a\le \eta,\qquad
  R(a): a\ge-\eta.
\]
At least one branch holds, and both may hold in the overlap
\([-\eta,\eta]\).  A node contract must prove that either branch gives an
acceptable output when its guard holds.  The constructor may choose the first
passing guard by bounded Boolean search.

\begin{lemma}[Overlapping-branch soundness]
If \(\abs{x-a}\le\eta\), \(L(a)\) implies \(x\le2\eta\), and \(R(a)\) implies
\(x\ge-2\eta\).  Moreover \(L(a)\vee R(a)\) is decidable and always true.
\end{lemma}

\begin{proof}
The first two claims are triangle inequalities.  For coverage, if
\(a\le\eta\) then \(L\) holds.  Otherwise the decidable rational inequality
gives \(a>\eta\ge-\eta\), so \(R\) holds.  This is a case split on an exact
rational Boolean.
\end{proof}

This suffices whenever the two branch outputs agree up to the allowed error
on the overlap.  It is the constructive analogue of a stable conditional.

\subsection{Exact decision from a gap}

\begin{theorem}[Gap-decision theorem]
Let \(x\) have a rational approximation \(a\) with
\(\abs{x-a}<\delta/4\), where \(\delta>0\).  Assume the finite evidence
\[
           x=0\quad\vee\quad \abs{x}\ge\delta.
\]
Then the rational test \(\abs a<\delta/2\) decides which alternative holds.
\end{theorem}

\begin{proof}
If \(\abs a<\delta/2\), the apart alternative would imply
\[
 \delta\le\abs x\le\abs{x-a}+\abs a
 <\delta/4+\delta/2<\delta,
\]
a contradiction; hence \(x=0\).  If \(\abs a\ge\delta/2\), then \(x=0\)
would give \(\abs a=\abs{a-x}<\delta/4\), a contradiction; hence
\(\abs x\ge\delta\).  The proof uses the supplied disjunction, not an
unrestricted law of excluded middle.
\end{proof}

The same proof decides among finitely many regions separated by a rational
margin.  It is the basis for exact pivot selection, inversion, square-free
decomposition with a discriminant bound, exact chart choice, and branch
collapse.

\subsection{Why the gap is not optional}

\begin{obstruction}[Uniform zero decision]
There is no total computable procedure that decides whether an arbitrary fast
Cauchy name of a nonnegative real represents zero.
\end{obstruction}

\begin{proof}
Let \(H(e,t)\) be the primitive-recursive Boolean saying that machine \(e\)
halts within \(t\) steps, and define
\[
  A(e,s)=\sum_{t<s}2^{-t-1}H(e,t).
\]
The omitted tail is at most \(2^{-s}\), so \(A(e,\cdot)\) is one fixed
primitive-recursive family of fast Cauchy names.  If the machine never halts,
the limit is zero.  If it first halts at stage \(h\), then \(H(e,t)=1\) for
all \(t\ge h\), and the limit is \(2^{-h}>0\).  A uniform zero decider would
therefore decide the halting problem.
\end{proof}

This metatheoretic statement does not prove that every first-order sentence
called ``root selection'' is unprovable in iPRA.  It identifies a particular
impossible interface: individually accessible exact choices from arbitrary
Cauchy names.  Approximate sets, finite covers, projective outputs, and
gap-conditioned selectors are different and often valid.

\subsection{Projective and quotient outputs}

Some apparent choices are artifacts of representation.  A rotation has two
unit-quaternion representatives \(q,-q\).  An eigenline has many nonzero
vectors and two unit vectors of opposite sign.  A polynomial's unordered
root multiset has many orderings.  The primary output should therefore be a
quotient-invariant certificate:
\[
  qq^\ast=1,\qquad R=R(q),
\]
or a rank-one projector, or a symmetric list code.  Normalization and sign
selection are secondary operations guarded by nonzero and orientation
evidence.

This move is not merely philosophical.  Quotient-invariant polynomial
identities are exact finite arithmetic and compose well in iPRA, while a
canonical representative may force a discontinuous branch.

\section{Bounded search, height certificates, and finiteness}

\subsection{Complete bounded enumeration}

Let \(P(d,x)\in\B\) be a primitive-recursive predicate, and let \(H(d)\) be a
primitive-recursive bound.

\begin{construction}[Bounded enumerator]
Scan \(x=0,\ldots,H(d)\).  Append \(x\) exactly when \(P(d,x)=1\).
Call the resulting coded list \(L(d,H)\).
\end{construction}

\begin{theorem}[Bounded enumeration]
The list \(L(d,H)\) is sound:
every listed member satisfies \(P\).  It is complete below \(H\):
if \(x\le H(d)\) and \(P(d,x)=1\), then \(x\) occurs in the list.
Both claims have quantifier-free induction proofs.
\end{theorem}

\begin{proof}
At stage \(t\), use the invariant:
the accumulator contains exactly the successful indices below \(t\), in
strictly increasing order.  ``Contains exactly'' is represented by two
Boolean folds: every listed entry is below \(t\) and satisfies \(P\), and
for every \(x<t\), \(P(d,x)=1\) implies a bounded membership test succeeds.
All ranges are bounded by \(t\) and compile to primitive-recursive folds.
The update at \(t\) preserves the invariant by a decidable test of \(P(d,t)\).
\end{proof}

\subsection{Height bounds and complete lists}

Suppose \(\mathsf{Sol}(d,x)\) is the exact coded solution predicate.

\begin{theorem}[Bound--list equivalence]
For a fixed decidable coding:
\begin{enumerate}[label=(\roman*)]
\item a verified height bound \(H\) produces a complete finite solution list
      by bounded enumeration;
\item a verified complete finite list produces a height bound, namely the
      maximum code height of its entries.
\end{enumerate}
The conversions are primitive recursive relative to the supplied
certificates.
\end{theorem}

\begin{proof}
Part (i) is the bounded enumeration theorem plus the certificate saying every
solution has height at most \(H\).  For part (ii), fold maximum over the list.
Completeness says every solution is equal to one of its finitely many
entries, so its height is at most that maximum.  Equality and list membership
are decidable in the numerical coding.
\end{proof}

The theorem does not turn an abstract statement
\(\forall d\,\exists H\,\forall x(\mathsf{Sol}(d,x)\to h(x)\le H)\)
into a primitive-recursive function \(H(d)\).  The production of \(H\) is a
separate theorem.

\subsection{Cardinality is not a cutoff}

\begin{theorem}[Singleton obstruction]
There is a uniformly primitive-recursive family of predicates \(P_e(t)\)
such that each set
\[
        S_e=\{t:P_e(t)=1\}
\]
has cardinality at most one, but there is no total computable function
\(B(e)\) bounding every member of \(S_e\).
\end{theorem}

\begin{proof}
Let \(P_e(t)=1\) exactly when machine \(e\) halts for the first time at stage
\(t\).  This is a primitive-recursive bounded simulation predicate, and
\(S_e\) is empty or a singleton.  If a total computable \(B(e)\) bounded its
member, simulate machine \(e\) for \(B(e)\) steps.  It halts eventually if and
only if it halts within that bound, so this would decide the halting problem.
\end{proof}

Consequently, a cardinality estimate for rational points cannot by itself
provide an enumeration algorithm or stopping criterion.  Pairwise separation
also does not help unless the ambient search region is bounded.

\subsection{Producer--verifier factorization}

For Diophantine and analytic number-theoretic packets, every claimed result
should be split into two modules.

\begin{description}[leftmargin=3.2em,style=nextline]
\item[Certificate producer.]
The deep mathematical theorem computes a height, cover, zero-free region,
exceptional modulus, or numerical constant.
\item[Certificate verifier.]
Finite arithmetic checks the certificate and derives the desired bounded
conclusion.
\end{description}

The verifier is often readily formalizable in iPRA even when the producer is
unknown, ineffective, or stronger than current mathematics.  This split is a
positive result: it preserves all finitary consequences of a future producer
and prevents the packet from disguising the missing theorem as an
implementation detail.

For general Faltings and Siegel targets, the framework supplies the exact
bounded enumerator and the equivalence between height bounds and complete
lists.  It does not supply the general height producer.  For prime
distribution, a certificate may contain explicit constants and verified
finite exceptional cases; the verifier proves the stated range once these are
present, but does not make an ineffective constant effective.

\section{Worked benchmark: the stable cubic \(y^3+y=x\)}

This section gives an end-to-end mathematical example suitable for the first
full object-language implementation of the framework.  It uses only exact
rational arithmetic, bounded bisection, and quantifier-free induction.

Fix a rational input \(x\) satisfying \(0\le x\le1\), and set
\[
       F_x(y)=y^3+y-x.
\]
We construct a regular name for the unique solution in \([0,1]\).

\subsection{Exact bisection}

Set \(\ell_0=0\) and \(u_0=1\).  Given \(\ell_n,u_n\), let
\[
      m_n=\frac{\ell_n+u_n}{2}.
\]
Since \(F_x(m_n)\) is rational, its sign is decidable exactly.  Define
\[
(\ell_{n+1},u_{n+1})=
\begin{cases}
 (m_n,u_n),&F_x(m_n)\le0,\\
 (\ell_n,m_n),&F_x(m_n)>0.
\end{cases}
\]

\begin{lemma}[Bisection invariant]
For every \(n\),
\[
  0\le\ell_n\le u_n\le1,\qquad
  F_x(\ell_n)\le0\le F_x(u_n),\qquad
  u_n-\ell_n=2^{-n}.
\]
\end{lemma}

\begin{proof}
At \(n=0\), \(F_x(0)=-x\le0\), while
\(F_x(1)=2-x\ge1\).  Suppose the invariant holds.  The midpoint lies between
the endpoints.  If \(F_x(m_n)\le0\), the left endpoint can be replaced by
\(m_n\); otherwise the right endpoint can be replaced by \(m_n\), and exact
rational order gives \(0<F_x(m_n)\).  In either case the sign bracket is
preserved and the width is halved.  The entire invariant is a conjunction of
rational cross-product inequalities and one dyadic-width equality, so its
induction matrix is quantifier-free.
\end{proof}

Let
\[
        c_n=\frac{\ell_n+u_n}{2}.
\]

\begin{lemma}[Regularity]
If \(m\ge n\), then
\[
          \abs{c_m-c_n}\le 2^{-n}.
\]
Consequently \(a(k)=c_{k+2}\) is a fast regular name.
\end{lemma}

\begin{proof}
The bisection intervals are nested.  A finite induction on \(m-n\) proves
\[
       [\ell_m,u_m]\subseteq[\ell_n,u_n].
\]
Both centers lie in the latter interval, whose width is \(2^{-n}\).  For the
shifted name,
\[
 \abs{a(j)-a(k)}
 \le 2^{-(k+2)}<2^{-k}
 \qquad(j\ge k).
\]
No completeness axiom is used in the internal statement; regularity is a
family of exact rational inequalities.
\end{proof}

\subsection{Residual bound}

For \(0\le r\le s\le1\),
\[
\begin{aligned}
  F_x(s)-F_x(r)
  &=(s-r)(s^2+sr+r^2+1),\\
  s-r
  &\le F_x(s)-F_x(r)\le4(s-r).
\end{aligned}
\]
The upper bound uses \(s^2,sr,r^2\le1\).

\begin{theorem}[Certified cubic name]
For every \(n\),
\[
       \abs{F_x(c_n)}\le 4\cdot2^{-n}.
\]
In particular, \(a(k)=c_{k+3}\) satisfies
\[
       \abs{a(k)^3+a(k)-x}\le2^{-k-1}.
\]
\end{theorem}

\begin{proof}
Monotonicity places \(F_x(c_n)\) between
\(F_x(\ell_n)\le0\) and \(F_x(u_n)\ge0\).  Hence
\[
 \abs{F_x(c_n)}
 \le F_x(u_n)-F_x(\ell_n)
 \le4(u_n-\ell_n)=4\cdot2^{-n}.
\]
All quantities are rational, and the factorization of the cubic difference
is a polynomial identity.
\end{proof}

\subsection{Quantitative uniqueness}

\begin{lemma}[Residual isolation]
For rational \(p,q\in[0,1]\),
\[
       \abs{p-q}\le \abs{F_x(p)-F_x(q)}
       \le \abs{F_x(p)}+\abs{F_x(q)}.
\]
Thus if both residuals are at most \(\eps\), then
\(\abs{p-q}\le2\eps\).
\end{lemma}

\begin{proof}
The factorization gives
\[
 \abs{F_x(p)-F_x(q)}
 =\abs{p-q}\,(p^2+pq+q^2+1).
\]
The second factor is at least \(1\), proving the first inequality.  The
second is the triangle inequality.
\end{proof}

It follows that any two regular names in \([0,1]\) whose residuals tend to
zero with explicit moduli are extensionally equal.  The proof is quantitative:
to establish distance at most \(2^{-k}\), request residuals at most
\(2^{-k-2}\) and approximations accurate enough that their rounding errors
consume the remaining half of the budget.

\subsection{The exact certificate}

A finite certificate for request \(k\) contains:
\[
  x,\quad n=k+3,\quad
  (\ell_0,u_0),\ldots,(\ell_n,u_n),
\]
together with the branch bit at each stage.  A verifier checks:
\begin{enumerate}[label=(\roman*)]
\item \(x\) and every endpoint are valid rational codes;
\item \(0\le x\le1\);
\item the initial interval is \([0,1]\);
\item each successor is the prescribed half selected by the recorded exact
      sign test;
\item every sign bracket and width equality holds;
\item the final center and residual bound are correct.
\end{enumerate}
The trace length is \(k+3\), a primitive-recursive function of the free
request.  The invariant is precisely the bisection invariant above.  This is
therefore a complete instance of the finite-trace reflection theorem.

\subsection{Extension to represented \(x\)}

If \(x\) is itself given by a regular name, exact midpoint signs are
unavailable near zero.  Use the overlapping-branch rule:
at stage \(n\), approximate \(F_x(m_n)\) within \(\eta_n\).  A confidently
negative value selects the right half, a confidently positive value selects
the left half, and an ambiguous value returns a short central interval whose
radius accounts for \(\eta_n\) using the lower slope bound \(1\).
The output is initially relational; the residual isolation lemma proves that
all accepted intervals are mutually close.  Reverse demand propagation picks
the input precision needed to make the final interval width and residual
smaller than \(2^{-k}\).  This demonstrates why the relational framework is
not a concession: it is the stable uniform form of bisection on represented
inputs.

\section{Application patterns and repairs of packet difficulties}

The framework is intended as a shared substrate.  This section states what it
does and does not resolve in the principal domains represented in the
repository.

\subsection{Algebraic equations and Puiseux expansions}

For a polynomial with exact rational or exact algebraic coefficient codes,
finite algebraic operations, resultants, discriminants, square-free
factorization certificates, and coefficient identities are decidable finite
data.  A root theorem should first output a finite list of rational disks with:
coverage of all roots, multiplicity counts, and a refinement relation between
precisions.  A separation certificate collapses an isolated disk to one root
name.

For arbitrary Cauchy-presented coefficients, exact Puiseux coefficient
selection is too strong.  The repository gives an explicit fixed-family
halting reduction using
\[
       P_a(X,Y)=Y^2-X^2-a.
\]
The repair is one of:
\begin{enumerate}[label=(\alph*)]
\item exact algebraic coefficient presentations with decidable zero and
      certified splitting;
\item a finite family of truncated approximate branches with residual and
      coverage bounds;
\item a theorem conditional on a simple-root or discriminant gap.
\end{enumerate}
The contract-DAG compiler handles the coefficient recursion once the branch
guard is legitimate; it does not create the missing exact equality oracle.

\subsection{Euler--Rodrigues, eigenvectors, and projective geometry}

Rodrigues' formula for a supplied axis and angle is a polynomial/trigonometric
error-contract problem.  Exact Euler-axis selection for an arbitrary
represented rotation is discontinuous at the identity and can encode LLPO.
The stable outputs are:
\begin{itemize}[leftmargin=1.6em]
\item a unit quaternion modulo \(q\sim-q\);
\item a projective axis or finite approximate axis set;
\item an exact axis under a supplied lower bound separating the rotation from
      the identity;
\item rational exact axes for rational proper orthogonal matrices.
\end{itemize}
The same projective-first discipline applies to eigenvectors.  Use a
projector, kernel certificate, or finite set of charts; normalize only under a
nonzero coordinate or singular-value gap.

Many Euclidean and projective geometry theorems become polynomial identities
after homogeneous coordinates are introduced.  Denominators should not be
cancelled silently.  Either clear them and retain nondegeneracy guards, or
work entirely with cross-products and determinants.  A finite incidence
certificate then fits directly into the contract calculus.

\subsection{Analysis and special functions}

A valid analysis node has:
\[
   \text{approximant},\quad
   \text{cutoff},\quad
   \text{tail/error verifier},\quad
   \text{precision demand}.
\]
For a series, the producer supplies \(N(k,d)\) and a rational proof that the
tail after \(N\) is at most \(2^{-k}\).  For an integral, it supplies a mesh
and local oscillation modulus.  For analytic continuation, it supplies a
finite chain of overlapping disks and derivative bounds.  The compiler adds
the budgets and verifies the finite chain.

This resolves a common formal difficulty in trigonometric, Gamma, zeta, and
residue packets: one never inducts on a completed real identity.  One proves
a rational inequality at arbitrary requested precision.  It does not, by
itself, prove a missing global zero-free region or manufacture an effective
constant from an ineffective theorem.

\subsection{Probability}

Finite probability laws are rational mass arrays with a Boolean check that
masses are nonnegative and sum to one.  Convolution, moments, and
characteristic-function samples are exact finite sums.  A constructive CLT
packet should separate:
\begin{enumerate}[label=(\roman*)]
\item the exact finite convolution engine;
\item moment and variance guards;
\item a finite-grid comparison or characteristic-function approximation;
\item the analytic smoothing and tail estimates with explicit constants.
\end{enumerate}
The first three are generic contract nodes.  The fourth is the substantive
analytic producer.

\subsection{Topology}

For Jordan-curve and invariance-of-dimension arguments, use finite grids,
rational boxes, adjacency graphs, parity or degree certificates, and explicit
separation margins.  A finite grid calculation at a few scales is evidence,
not a proof for every scale.  The reusable theorem must show that the
quantifier-free grid invariant is preserved under refinement.

Degenerate boundary cases should be handled by overlap: several adjacent
cells may be valid.  A positive distance from the boundary collapses the
finite relation to one cell.  This is the topological instance of the same
branching theorem used for roots and axes.

\subsection{Diophantine geometry}

Runge-type arguments often do produce explicit polynomial identities and
height inequalities.  These are ideal evidence-indexed nodes: verify the
identity, derive the bound, then enumerate.  The general Siegel and Faltings
packets correctly identify the missing general height producer.  The
framework contributes:
\begin{itemize}[leftmargin=1.6em]
\item exact encodings of rational and integral points;
\item complete bounded enumeration;
\item conversion between a verified bound and a complete list;
\item finite-cover and finite-algebra verifiers;
\item a precise interface for any future height or covering theorem.
\end{itemize}
A cardinality theorem, even a strong one, cannot be plugged into the height
input slot.

\subsection{Analytic number theory}

Prime-distribution statements must expose all dependencies:
modulus range, exceptional characters, zero-free constants, conductor cutoff,
density estimate, and finite small cases.  A certificate verifier can check
the finite exceptional list and derive a claimed numerical range from
explicit inequalities.  The quantifier order is crucial:
\[
  \exists C\;\forall x\ge X(C)\;\Phi(C,x)
\]
is useful only when \(C\) is produced or supplied.  A proof of mere existence
in a stronger classical theory does not automatically yield a PR constant in
iPRA.  The framework therefore records ``analytic constant producer'' as a
first-class dependency rather than burying it in a phrase such as ``take
\(x\) sufficiently large.''

\subsection{Summary matrix}

\begin{longtable}{p{.22\linewidth}p{.34\linewidth}p{.34\linewidth}}
\toprule
\textbf{Domain} & \textbf{Uniform output supported} &
\textbf{Additional evidence for exact choice}\\
\midrule
\endhead
Polynomial roots & finite disk cover, residuals, multiplicities &
root separation or exact algebraic splitting\\
Puiseux & finite approximate truncations &
exact coefficient field and branch isolation\\
Rotations & quaternion modulo sign, approximate axis set &
distance from identity / nonzero vector part\\
Eigenvectors & projective line or projector &
spectral gap and normalization chart\\
Limits and integrals & rational approximants with moduli &
none beyond the quantitative hypotheses\\
Topology & finite grid/parity relation &
boundary separation for unique cell/branch\\
Diophantine points & complete list below supplied height &
height or covering certificate producer\\
Prime distribution & verifier for explicit constants/ranges &
effective constant and exceptional-set producer\\
\bottomrule
\end{longtable}

\section{Lean formalization architecture}

The archive contains two Lean blueprints.  They are deliberately separated
into a semantic algorithm layer and a deeply embedded object-theory layer.

\subsection{Layer 1: executable semantic algorithms}

The file \code{lean/FiniteContractSemantic.lean} defines:
\begin{itemize}[leftmargin=1.6em]
\item bounded Boolean folds and first-witness search;
\item finite trace execution;
\item reverse demand propagation on a topologically ordered finite graph;
\item fuelled backward viability and coherent finite paths;
\item bounded enumeration;
\item rational cubic bisection specifications.
\end{itemize}
Its intended theorems are ordinary Lean statements about \code{Nat},
\code{List}, and \code{Bool}.  They serve three purposes: test the definitions,
clarify invariants, and provide external semantics lemmas.  They are not the
iPRA proof.

\subsection{Layer 2: actual ProveIt interfaces}

The object blueprint imports:
\begin{lstlisting}[language={}]
import ConciseIPRA.PR
import ConciseIPRA.NativeBiInterpretation
\end{lstlisting}
and works in \code{LeanProofs.ConciseIPRA}.  The relevant existing declarations
are:
\begin{lstlisting}[language={}]
PR.zero   PR.succ   PR.proj   PR.comp   PR.prec
PTm.var   PTm.app
PRAAx.proj  PRAAx.comp  PRAAx.precZero  PRAAx.precSucc
PRAAx.succNeZero  PRAAx.succInj  PRAAx.ind
PRA.proj  PRA.comp  PRA.precZero  PRA.precSucc
PRA.qfInd
pra_iff_native
\end{lstlisting}

The first implementation milestone is a small object-level arithmetic
library: numerals, addition, multiplication, truncated subtraction, order
Booleans, pair/list coding, Boolean folds, and table lookup.  Each function is
a concrete \code{PR} term, and each defining equation is a
\(\Prf\,\PRAAx\) theorem composed from \code{PRA.comp},
\code{PRA.precZero}, and \code{PRA.precSucc}.

\subsection{Object trace structure}

A practical record is:
\begin{lstlisting}[language={}]
structure PRTraceSystem (p : Nat) where
  init    : PR p
  step    : PR (p + 2)   -- stage, state, parameters
  guard   : PR (p + 1)   -- state, parameters -> Boolean
  finish  : PR (p + 1)
  horizon : PR p
\end{lstlisting}
The exact arities can be adjusted to the repository's convention that the
recursion argument comes first.  The invariant is not a Lean predicate on
natural numbers; it is an \code{Fm PTm} whose atomic components assert that
the Boolean \code{guard} term equals \(1\).

The core target has the schematic shape:
\begin{lstlisting}[language={}]
theorem PRA.traceReflection
    (S : PRTraceSystem p)
    (hQF : (traceInv S).QF)
    (hBase : Gamma |-[(PRAAx)] baseFormula S)
    (hStep : stepContext S Gamma |-[(PRAAx)] stepFormula S)
    (hFinish : finishContext S Gamma |-[(PRAAx)] outputFormula S) :
    Gamma |-[(PRAAx)] finalFormula S := by
  -- instantiate PRA.qfInd with the trace invariant;
  -- use S.init/S.step defining equations;
  -- compose hBase, hStep, hFinish.
\end{lstlisting}
The actual file uses ASCII comments rather than claiming the notation above
already elaborates.

\subsection{Compilation milestones}

\begin{enumerate}[label=\textbf{M\arabic*.},leftmargin=3em]
\item \textbf{Boolean folds.}
      Define \code{prAllBelow}, \code{prAnyBelow}, and \code{prFirst}.
      Prove their zero/successor equations and object-level correctness.
\item \textbf{Coded tables.}
      Define pair, projections, list length, lookup, append, and update as
      \code{PR} terms.  Prove the lookup-after-update equations needed by the
      trace invariant.
\item \textbf{Generic trace.}
      Encode \code{run} by primitive recursion.  Prove its zero and successor
      equations using \code{PRA.precZero}/\code{PRA.precSucc}.
\item \textbf{Trace reflection.}
      Express the invariant as an \code{Fm PTm}; prove quantifier-freeness;
      invoke \code{PRA.qfInd}.
\item \textbf{Reverse demand compiler.}
      For a fixed graph descriptor generated at Lean elaboration time, build
      the concrete \code{PR} reverse fold and prove edge adequacy.
\item \textbf{Node combinators.}
      Implement identity, deterministic map, product, guarded branch,
      bounded search, relational refinement, and sequential composition.
\item \textbf{DAG theorem.}
      Generate the global invariant from the node list and compose the local
      \code{Prf} objects.
\item \textbf{Cubic benchmark.}
      Instantiate the generic trace with rational bisection and export the
      all-precision residual theorem.
\item \textbf{Native transfer.}
      Only after the \PRAAx\ theorem is complete, apply
      \code{pra\_iff\_native} or \code{NFm.ofPRA} to obtain a theorem in the
      native scalar core.
\end{enumerate}

\subsection{Reflection versus theorem generation}

There are two complementary strategies.

\paragraph{Concrete certificate reflection.}
For one closed input and one requested precision, a small verified checker
can reduce the certificate to computation.  A checker soundness theorem then
converts successful evaluation into a proof.

\paragraph{Uniform theorem generation.}
For all inputs, finite computation is not enough.  The generated proof must
contain the base case, step case, and quantifier-free induction.  A
meta-level compiler may construct these proof terms from a graph descriptor,
but the resulting terms are checked by Lean's kernel against the \PRAAx\
rules.

The second strategy is the long-term goal.  It avoids trusting a large tactic:
the compiler may be untrusted because it emits explicit \code{Prf} objects.

\subsection{Formalization status}

The Lean sources supplied here are a detailed blueprint and semantic
prototype.  Lean 4.32.0 was not available in the artifact-building runtime,
so the files were not compiler-checked.  The archive marks every unfinished
object theorem with \code{sorry} and never treats it as established.
The mathematical proofs in this report identify the required invariants and
the exact existing constructors needed to remove those placeholders.

\section{A conversion workflow for existing proof packets}

The repository's packets differ greatly in maturity.  The following workflow
turns each into a comparable, auditable artifact.

\subsection{Step 1: state the constructive theorem}

Replace a classical slogan by an explicit input/output contract.  State:
the numerical input code, well-formedness test, request, output relation,
error metric, and every extra evidence item.  Separate unconditional
approximate output from gap-conditioned exact output.

\subsection{Step 2: classify every dependency}

Each dependency receives one of five labels:
\[
\begin{array}{ll}
\mathsf{PR}&\text{fixed primitive-recursive computation},\\
\mathsf{LOCAL}&\text{proved finite arithmetic contract},\\
\mathsf{CERT}&\text{supplied finite certificate with verifier},\\
\mathsf{PRODUCER}&\text{theorem that constructs such a certificate},\\
\mathsf{BARRIER}&\text{interface ruled out or known to require stronger data}.
\end{array}
\]
This prevents a missing producer from being mistaken for a coding problem.

\subsection{Step 3: build the contract graph}

Make each substantive lemma a node.  Edges record both mathematical
dependence and precision demand.  Common nodes---rational arithmetic,
polynomial evaluation, finite sums, bounded search, interval refinement,
finite nets---should point to shared libraries rather than being copied.

\subsection{Step 4: totalize operations}

Every operation must be total on all number codes.  Division, normalization,
chart inversion, and array access return a harmless default outside their
domain.  A Boolean guard is carried in the invariant and proves that the
default branch is never used in a valid run.

\subsection{Step 5: compute reverse demands}

Starting from the theorem's requested output precision, compute every
upstream precision, grid scale, cutoff, and iteration count.  Record these
functions explicitly.  Phrases such as ``take \(N\) large enough'' are not
accepted as a contract.

\subsection{Step 6: isolate discontinuities}

For every exact branch ask:
\begin{enumerate}[label=(\alph*)]
\item Can the output be quotient-invariant or relational?
\item Can overlapping guards make the approximate theorem uniform?
\item Which positive gap would make the exact branch decidable?
\item Is the requested selector contradicted by a known computability or
      constructive-logic reduction?
\end{enumerate}

\subsection{Step 7: define the finite trace invariant}

The invariant must mention only:
loop counters, coded states, Boolean verifiers, bounds, and finite folds.
If it contains ``for every precision'' or quantification over a function, the
proof has not yet been compiled to the correct level.

\subsection{Step 8: prove semantic algorithms}

Implement exact reference code and ordinary Lean definitions.  Test malformed
inputs and boundary cases.  These artifacts discover mistakes, but their
status is explicitly ``semantic'', not ``embedded''.

\subsection{Step 9: construct object proofs}

Translate each fixed algorithm to \code{PR}; prove defining equations;
construct local \code{Prf PRAAx} contracts; use the generic trace theorem.
Record the remaining \code{sorry} count and axiom audit.  A file containing
only proposition definitions belongs in the ``targets'' category.

\subsection{Step 10: transport and interpret}

After the \PRAAx\ proof exists, transport it to the native core if desired.
Then prove an external interpretation theorem connecting the codes to the
ordinary theorem.  This last theorem may use a richer Lean mathematics
library without changing what was proved internally.

\subsection{Recommended machine-readable status}

Every packet should contain a ledger with fields:
\begin{lstlisting}[language={}]
{
  "claim": "...",
  "input_presentation": "...",
  "output_relation": "...",
  "producer_status": "proved | conditional | missing | obstructed",
  "verifier_status": "human | semantic_lean | object_prf",
  "precision_transformer": "...",
  "barriers": ["..."],
  "source_commit": "..."
}
\end{lstlisting}

The status values are monotone.  A semantic Lean theorem does not upgrade
itself to \code{object\_prf}; compiled definitions do not upgrade a target to
a theorem; and a finite example does not upgrade a universal claim.

\subsection{Suggested repository library}

A reusable library could be organized as:
\begin{lstlisting}[language={}]
Logic/PrimitiveRecursiveArithmetic/CFP/
  Coding.lean
  BoolFold.lean
  Rational.lean
  FiniteTrace.lean
  Contract.lean
  Demand.lean
  Relational.lean
  BoundedSearch.lean
  FiniteNet.lean
  ObjectCompiler.lean
  Examples/Cubic.lean
\end{lstlisting}
The research packets would import this library and provide only their
domain-specific nodes and producers.  This is the principal engineering gain
of the framework: the difficult proof-theoretic plumbing is done once.

\section{Status, limitations, and further research}

\subsection{What is proved in this report}

The human-readable arguments establish:
\begin{itemize}[leftmargin=1.6em]
\item Boolean compilation of bounded quantifiers and bounded witness search;
\item finite-trace reflection from base and step contracts;
\item the contract-DAG composition theorem with reverse demands;
\item finite-horizon coherent selection for relational outputs;
\item collapse to deterministic output under an isolation certificate;
\item overlapping-branch soundness and exact gap decision;
\item backward precision budgeting;
\item bounded enumeration and the bound--list equivalence;
\item the impossibility of deriving a cutoff from a uniform cardinality-one
      bound;
\item an exact bisection certificate and quantitative uniqueness theorem for
      \(y^3+y=x\) on \(0\le x\le1\).
\end{itemize}

The proofs are finitary and explicitly identify quantifier-free invariants.
They are designed to be formalized with the inspected \code{ConciseIPRA}
interfaces.

\subsection{What is not claimed}

No complete generic \(\Prf\,\PRAAx\) implementation is included.  The Lean
files are blueprints and semantic prototypes, not compiler-checked object
proofs.  The report does not prove:
\begin{itemize}[leftmargin=1.6em]
\item general effective Siegel or Faltings;
\item a primitive-recursive generator of Mordell--Weil bases for arbitrary
      elliptic curves;
\item effective versions of every prime-distribution estimate in the packet
      catalog;
\item uniform exact Puiseux branches from arbitrary Cauchy coefficients;
\item exact unit-axis selection for arbitrary represented rotations;
\item an infinite-path principle for arbitrary finitely branching trees.
\end{itemize}

In each case, the framework gives either a conditional verifier theorem, a
weaker stable output, or a barrier explaining why the stronger interface
requires additional data.

\subsection{Priority formalization targets}

\begin{enumerate}[label=\arabic*.]
\item \textbf{Object Boolean folds.}  They unlock bounded quantifiers,
      finite lists, search, and graph invariants.
\item \textbf{Generic finite trace.}  This should be the first substantial
      reusable \code{Prf PRAAx} theorem.
\item \textbf{Cubic benchmark.}  It exercises rational codes, branching,
      recursion, inequalities, and all-precision introduction without deep
      external mathematics.
\item \textbf{Relational finite roots.}  A polynomial disk-cover benchmark
      tests compatibility and multivalued outputs.
\item \textbf{One existing packet.}  Basel-type finite inequalities or a
      Runge bounded-enumeration case would test end-to-end integration.
\end{enumerate}

\subsection{Research questions}

\begin{enumerate}[label=\textbf{Q\arabic*.},leftmargin=3em]
\item Can the local contract language be chosen so that quantifier-freeness
      and proof composition are guaranteed by construction?
\item Can an untrusted Lean metaprogram compile a contract graph to compact
      \code{Prf PRAAx} terms whose size is polynomial in the graph and local
      proofs?
\item Which classical single-valued theorems admit uniform finite relational
      reformulations, and what quantitative isolation data collapses them?
\item Can proof mining systematically extract the reverse demand maps from
      existing constructive Lean proofs?
\item Which ineffective analytic theorems have independently verifiable
      finite certificates even when no known uniform producer exists?
\item Can a common finite-cover interface serve topology, complex analysis,
      algebraic root isolation, and Diophantine subspace arguments?
\item What is the smallest native-core proof checker needed for the generated
      contract proofs, measured after maximal gzip compression?
\end{enumerate}

\subsection{Conclusion}

The right general framework for mathematics in iPRA is not a library of
completed infinite objects.  It is a calculus of finite observations,
explicit demands, guarded operations, relational outputs, and bounded proof
traces.  The contract-DAG theorem gives these ingredients a compositional
form.  It makes successful packet proofs reusable and failed packet proofs
diagnosable: every gap appears as a missing producer, missing verifier,
missing quantitative margin, or impossible selector.  That is enough to move
from isolated outlines toward a genuine formal library.


\appendix
\section{Repository source audit}

The repository was inspected read-only at
\[
 \code{b88e843c31b50dfe461bfbd0380166ca422cd3ab}.
\]
The following paths were central to this synthesis.

\begin{longtable}{p{.39\linewidth}p{.52\linewidth}}
\toprule
\textbf{Path} & \textbf{Role in this report}\\
\midrule
\endhead
\path{Logic/PrimitiveRecursiveArithmetic/README.md} &
Defines the checked/research split, reports fifty-eight proof packets, and
states that no packet currently contains a completed \code{Prf PRAAx}
derivation of its main theorem.\\

\path{ConciseIPRA/Lean/ConciseIPRA/PR.lean} &
Defines \code{PR}, \code{PTm}, \code{PRAAx}, the defining-equation theorems,
and \code{PRA.qfInd}.  These are the exact object-language interfaces used in
the Lean plan.\\

\path{ConciseIPRA/Lean/ConciseIPRA/NativeBiInterpretation.lean} &
Provides \code{pra\_iff\_native}, \code{native\_iff\_pra}, and round-trip
theorems.  It confirms that transport follows, rather than replaces, an
object proof.\\

\path{docs/proofs/foundations/ipra_framework_2026-11-15/} &
Certificate-and-modulus framework: finite presentations, quantitative names,
apartness, bounded search, finite descent, and a clear semantic/object proof
boundary.\\

\path{docs/proofs/foundations/ipra_embedding_framework/} &
Finite observations and demand, local analytic branches, finite windows,
hyperelliptic section counts, and finite algebra certificates.\\

\path{docs/proofs/foundations/ipra_evidence_framework/} &
Finite relational-output and compatibility interface; backward finite
dynamic programming for witness amalgamation.\\

\path{docs/proofs/foundations/ipra_finite_observation/} &
Finite observation compiler, fixed-evaluator factorization, finite-net
minimizers, and superelliptic section bounds, with explicit boundary
statements.\\

\path{docs/proofs/foundations/ipra_framework/} &
Finite contracts, stable windows, coherent approximate witnesses, analytic
branches, and integral power-gap certificates.\\

\path{docs/proofs/foundations/ipra_uniform_certificates/} &
Finite-net lifting, sign-bracket root names, and precision-contract
composition.\\

\path{docs/proofs/foundations/ipra_embedding_framework_2026-11-15/} &
The current 40-page Certified Finitary Presentation baseline.  It develops
finite traces, robust branching, finite nets, height/list conversion, and
barrier theorems.  The present report adds the explicit contract-DAG
compiler and packet-conversion architecture.\\

\path{docs/proofs/algebra/puiseux_ipra/evidence/scope_obstruction.md} &
Explicit fixed primitive-recursive halting-family obstruction to uniform
exact Puiseux coefficient selection from arbitrary Cauchy coefficients.\\

\path{docs/proofs/geometry/euler_rodrigues_ipra/README.md} &
Records the LLPO obstruction to unrestricted exact axis selection and the
positive rational, gap-conditioned, approximate, and quaternion variants.\\

\path{docs/proofs/diophantine-number-theory/faltings_ipra/README.md} &
Separates cardinality from height, explains the effective Mordell consequence
of a positive uniform height theorem, and supplies constructive special
cases only.\\

\path{docs/proofs/diophantine-number-theory/siegel_ipra/README.txt} &
Supplies a constructive Runge subcase and identifies the unresolved general
height/list and cardinality targets.\\
\bottomrule
\end{longtable}

\subsection{Method of use}

The foundational packet results are credited as prior components.  This
report does not claim that the packet statements have been independently
formalized or verified.  It uses their stated contracts and obstacle
analyses as design input, cross-checks them against the actual object-language
API, and gives new composition theorems at the human-proof level.

The Git history changed during the audit: the Certified Finitary Presentation
packet was placed on \code{main} while the work was in progress.  The final
snapshot above includes that merge.  This is why the report treats that packet
as baseline rather than presenting a parallel reinvention.

\section{Exact Lean API map}

This appendix records the minimum interface needed to complete the formal
development.

\subsection{\code{PR} and terms}

The repository defines:
\begin{lstlisting}[language={}]
inductive PR : Nat -> Type
  | zero : PR 0
  | succ : PR 1
  | proj {k : Nat} (i : Fin k) : PR k
  | comp {m k : Nat} (f : PR m) (gs : Fin m -> PR k) : PR k
  | prec {k : Nat} (f : PR k) (g : PR (k + 2)) : PR (k + 1)

inductive PTm (n : Nat) : Type
  | var (i : Fin n) : PTm n
  | app {k : Nat} (f : PR k) (ts : Fin k -> PTm n) : PTm n
\end{lstlisting}
The recursion argument of \code{prec} comes first:
\[
\begin{aligned}
 \mathsf{prec}(f,g)(0,\vec x)&=f(\vec x),\\
 \mathsf{prec}(f,g)(S y,\vec x)
   &=g(y,\mathsf{prec}(f,g)(y,\vec x),\vec x).
\end{aligned}
\]

\subsection{\PRAAx}

The axiom family has constructors for projections, composition, zero and
successor recursion equations, \(Sx\ne0\), successor injectivity, and
quantifier-free induction:
\begin{lstlisting}[language={}]
PRAAx.proj
PRAAx.comp
PRAAx.precZero
PRAAx.precSucc
PRAAx.succNeZero
PRAAx.succInj
PRAAx.ind (phi : Fm PTm (n + 1)) (h : phi.QF)
\end{lstlisting}
The namespace \code{PRA} provides substitution-ready theorem wrappers:
\begin{lstlisting}[language={}]
PRA.proj
PRA.comp
PRA.precZero
PRA.precSucc
PRA.succNeZero
PRA.qfInd
\end{lstlisting}

The signature of \code{PRA.qfInd} is conceptually:
\[
\frac{
 \Gamma\vdash\varphi(0)
 \qquad
 \varphi(y),\Gamma^\uparrow\vdash\varphi(Sy)
 \qquad \varphi\text{ quantifier-free}
}{
 \Gamma\vdash\forall y\,\varphi(y)
}.
\]
The formalization should use this theorem directly rather than rebuild
induction from the raw axiom.

\subsection{Proof composition}

The deduction library supplies hypothesis, axiom, implication introduction
and elimination, conjunction operations, universal introduction and
elimination, and substitution.  Local node contracts should be packaged as
ordinary Lean values of the relevant \code{Prf} type.  The graph compiler is
then a meta-level recursive function that combines those values.

A useful first representation is:
\begin{lstlisting}[language={}]
structure LocalContract (n : Nat) where
  premise : Fm PTm n
  result  : Fm PTm n
  premiseQF : premise.QF
  resultQF  : result.QF
  proof : [premise] |-[(PRAAx)] result
\end{lstlisting}
The actual context notation and universe parameters should follow
\code{ConciseIPRA.Deduction}; the blueprint file keeps these details explicit.

\subsection{Native transfer}

The repository theorem
\begin{lstlisting}[language={}]
pra_iff_native :
  (Gamma |-[(PRAAx)] phi) <->
  NPrf (Gamma.map NFm.ofPRA) (NFm.ofPRA phi)
\end{lstlisting}
is the final bridge.  It should appear at the end of the dependency graph.
Using it earlier does not reduce the obligation to construct the left-hand
derivation.

\subsection{Recommended theorem names}

\begin{longtable}{p{.34\linewidth}p{.56\linewidth}}
\toprule
\textbf{Name} & \textbf{Statement}\\
\midrule
\endhead
\code{PRA.allBelow\_correct} &
Boolean fold equals \(1\) exactly when all bounded tests succeed.\\
\code{PRA.firstTrue\_correct} &
Bounded witness routine returns a valid least witness or the sentinel.\\
\code{PRA.tableRun\_zero}, \code{tableRun\_succ} &
Defining equations of the coded trace.\\
\code{PRA.traceReflection} &
Base/step/final local proofs imply the final trace verifier.\\
\code{PRA.reverseDemand\_adequate} &
Every upstream request dominates each outgoing demand.\\
\code{PRA.relationalPath} &
Finite viability yields a compatible path through a requested horizon.\\
\code{PRA.gapDecision} &
A fine approximation plus zero-or-gap evidence decides the alternative.\\
\code{PRA.contractDAGReflection} &
Local node proof objects compose to the global sink contract.\\
\code{PRA.cubicBisection} &
The cubic trace returns a regular residual name with quantitative uniqueness.\\
\bottomrule
\end{longtable}

\section{Reference algorithms}

The archive's Python companion uses exact integers and
\code{fractions.Fraction}.  The algorithms are included for reproducibility
and for comparison with future Lean definitions.

\subsection{Bounded witness}

\begin{lstlisting}[language=Python]
def first_true(bound, pred):
    for i in range(bound):
        if pred(i):
            return i
    return bound
\end{lstlisting}
The sentinel is the bound itself.  The object implementation uses primitive
recursion rather than a Python loop, but has the same invariant.

\subsection{Reverse demands}

For topologically numbered nodes:
\begin{lstlisting}[language=Python]
req = external_requests[:]
for v in reversed(range(number_of_nodes)):
    for u, demand in parents[v]:
        req[u] = max(req[u], demand(req[v]))
\end{lstlisting}
In a formal packet, the graph and demand functions are fixed data.  The
Python function accepts them as ordinary objects only for testing.

\subsection{Backward viability}

\begin{lstlisting}[language=Python]
viable[K] = set(candidates[K])
for k in reversed(range(K)):
    viable[k] = {
        c for c in candidates[k]
        if any(compat(k, c, d) and d in viable[k + 1]
               for d in candidates[k + 1])
    }
\end{lstlisting}
A forward pass chooses the least compatible viable successor.  The formal
implementation represents each set by a Boolean table bounded by the
candidate-list length.

\subsection{Cubic bisection}

\begin{lstlisting}[language=Python]
lo, hi = Fraction(0), Fraction(1)
for _ in range(steps):
    mid = (lo + hi) / 2
    if mid**3 + mid - x <= 0:
        lo = mid
    else:
        hi = mid
return lo, hi
\end{lstlisting}
The tests verify the sign bracket, exact dyadic width, residual bound,
regularity across horizons, and quantitative uniqueness on a rational grid.

\subsection{Validation interpretation}

Passing tests establishes that the executable definitions behave as intended
on the tested cases.  It does not establish the universal theorems and does
not construct a \code{Prf PRAAx}.  The value of the reference code is narrower
and still important: it detects encoding mistakes, wrong precision shifts,
dead relational branches, invalid sentinels, and boundary-condition errors
before the object proof is attempted.


\begin{thebibliography}{99}

\bibitem{BishopBridges}
E.~Bishop and D.~Bridges,
\emph{Constructive Analysis},
Springer, 1985.

\bibitem{BridgesRichman}
D.~Bridges and F.~Richman,
\emph{Varieties of Constructive Mathematics},
Cambridge University Press, 1987.

\bibitem{TroelstraVanDalen}
A.~S.~Troelstra and D.~van Dalen,
\emph{Constructivism in Mathematics}, Volumes I--II,
North-Holland, 1988.

\bibitem{Weihrauch}
K.~Weihrauch,
\emph{Computable Analysis},
Springer, 2000.

\bibitem{Kohlenbach}
U.~Kohlenbach,
\emph{Applied Proof Theory: Proof Interpretations and Their Use in Mathematics},
Springer, 2008.

\bibitem{HajekPudlak}
P.~Hájek and P.~Pudlák,
\emph{Metamathematics of First-Order Arithmetic},
Springer, 1993.

\bibitem{Lean}
L.~de Moura, S.~Kong, J.~Avigad, F.~van Doorn, and J.~von Raumer,
The Lean theorem prover,
in \emph{Automated Deduction---CADE-25}, 2015.

\bibitem{ProveIt}
V.~Reshetnikov,
\emph{ProveIt}, repository snapshot
\code{b88e843c31b50dfe461bfbd0380166ca422cd3ab},
15 November 2026 (Pacific time).

\end{thebibliography}

\end{document}
